% ============================================================================
% WORKBOOK (MODIFIED) — Section 3.1: Integers and Their Opposites
% CCSS 7.NS.A.1 (VA: SOL rational number ordering emphasis)
% Practice-focused reinforcement with ordering/comparing rational numbers
% ============================================================================
\section{Integers and Their Opposites}
\topicTitle{Integers and Their Opposites}

\begin{quickReview}{Integers, Opposites, and Comparing Rational Numbers}
\textbf{Integers} are whole numbers and their negatives: $\ldots, -3, -2, -1, 0, 1, 2, 3, \ldots$

\begin{itemize}[leftmargin=*]
    \item Every integer has an \textbf{opposite}: the number the same distance from $0$ on the other side.
    \item \textbf{Absolute value} $|n|$ is the distance from $0$. Always $\geq 0$.
    \item To \textbf{compare} rational numbers in different forms (fractions, decimals, percents), convert them to the same form first.
    \item On a number line, numbers increase from left to right.
\end{itemize}

\medskip
\textbf{Example:} Compare $-\frac{3}{4}$ and $-0.6$. Convert: $-\frac{3}{4} = -0.75$. Since $-0.75 < -0.6$, we have $-\frac{3}{4} < -0.6$.
\end{quickReview}

% ── Warm-Up ──────────────────────────────────────────────────────────────────
\begin{practiceBox}[funGreen]{\faSun~Warm-Up}

\practiceHeader[funGreen]{Find the Opposite and Absolute Value}

\begin{multicols}{2}
\prob What is the opposite of $9$? \answerBlank[2cm]
\answerExplain{$-9$}{The opposite of $9$ is $-9$. Both are $9$ units from $0$.}

\prob What is the opposite of $-15$? \answerBlank[2cm]
\answerExplain{$15$}{The opposite of a negative is positive: $-(-15) = 15$.}

\prob Find $|{-22}|$. \answerBlank[2cm]
\answerExplain{$22$}{$-22$ is $22$ units from $0$, so $|{-22}| = 22$.}

\prob Find $|17|$. \answerBlank[2cm]
\answerExplain{$17$}{$17$ is $17$ units from $0$, so $|17| = 17$.}

\prob What is $13 + (-13)$? \answerBlank[2cm]
\answerExplain{$0$}{A number plus its opposite always equals $0$.}

\prob What is the opposite of $0$? \answerBlank[2cm]
\answerExplain{$0$}{Zero is its own opposite.}
\end{multicols}
\end{practiceBox}

% ── Practice A: Comparing Rational Numbers ───────────────────────────────────
\begin{practiceBox}{Comparing Rational Numbers}

\practiceHeader{Use $<$, $>$, or $=$ to compare. Convert to the same form.}

\prob $-\dfrac{2}{5}$ \answerBlank[1.5cm] $-0.3$
\answerExplain{$<$}{$-\frac{2}{5} = -0.4$. Since $-0.4 < -0.3$, the fraction is less.}

\prob $0.75$ \answerBlank[1.5cm] $\dfrac{4}{5}$
\answerExplain{$<$}{$\frac{4}{5} = 0.8$. Since $0.75 < 0.8$, the decimal is less.}

\prob $-\dfrac{1}{3}$ \answerBlank[1.5cm] $-\dfrac{2}{6}$
\answerExplain{$=$}{$-\frac{2}{6} = -\frac{1}{3}$. They are equal.}

\prob $-1.25$ \answerBlank[1.5cm] $-\dfrac{5}{4}$
\answerExplain{$=$}{$-\frac{5}{4} = -1.25$. They are equal.}

\prob $\dfrac{3}{8}$ \answerBlank[1.5cm] $0.4$
\answerExplain{$<$}{$\frac{3}{8} = 0.375$. Since $0.375 < 0.4$, the fraction is less.}

\prob $-\dfrac{7}{10}$ \answerBlank[1.5cm] $-\dfrac{3}{5}$
\answerExplain{$<$}{$-\frac{3}{5} = -0.6$ and $-\frac{7}{10} = -0.7$. Since $-0.7 < -0.6$, we have $-\frac{7}{10} < -\frac{3}{5}$.}

\end{practiceBox}

% ── Practice B: Ordering Rational Numbers ────────────────────────────────────
\begin{practiceBox}[funTeal]{Ordering Rational Numbers}

\practiceHeader[funTeal]{Order each set from least to greatest. Convert to decimals.}

\prob $\dfrac{1}{2}, \; -0.75, \; -\dfrac{1}{4}, \; 0.3, \; -1$ \answerBlank[6cm]
\answerExplain{$-1, \; -0.75, \; -\frac{1}{4}, \; 0.3, \; \frac{1}{2}$}{Decimals: $0.5, -0.75, -0.25, 0.3, -1$. From least: $-1 < -0.75 < -0.25 < 0.3 < 0.5$.}

\prob $-\dfrac{3}{4}, \; 0.8, \; -0.5, \; \dfrac{2}{3}, \; -\dfrac{1}{8}$ \answerBlank[6cm]
\answerExplain{$-\frac{3}{4}, \; -0.5, \; -\frac{1}{8}, \; \frac{2}{3}, \; 0.8$}{Decimals: $-0.75, 0.8, -0.5, 0.667, -0.125$. Order: $-0.75 < -0.5 < -0.125 < 0.667 < 0.8$.}

\prob $1.5, \; -\dfrac{7}{4}, \; \dfrac{3}{2}, \; -1.8, \; 0$ \answerBlank[6cm]
\answerExplain{$-1.8, \; -\frac{7}{4}, \; 0, \; \frac{3}{2}, \; 1.5$}{Decimals: $1.5, -1.75, 1.5, -1.8, 0$. Note $\frac{3}{2} = 1.5$. Order: $-1.8 < -1.75 < 0 < 1.5 = 1.5$.}

\prob $-\dfrac{5}{6}, \; -0.9, \; -\dfrac{4}{5}, \; -\dfrac{7}{8}$ \answerBlank[6cm]
\answerExplain{$-0.9, \; -\frac{7}{8}, \; -\frac{5}{6}, \; -\frac{4}{5}$}{Decimals: $-0.833, -0.9, -0.8, -0.875$. Order: $-0.9 < -0.875 < -0.833 < -0.8$.}

\end{practiceBox}

% ── Practice C: Mixed Form Comparison ────────────────────────────────────────
\begin{practiceBox}[funOrange]{Mixed Forms --- Fractions, Decimals, and Percents}

\practiceHeader[funOrange]{Compare by converting to the same form.}

\prob Which is greater: $\dfrac{3}{5}$ or $55\%$? \answerBlank[3cm]
\answerExplain{$\frac{3}{5}$}{$\frac{3}{5} = 60\%$. Since $60\% > 55\%$, the fraction is greater.}

\prob Order from least to greatest: $40\%, \; \dfrac{3}{8}, \; 0.39$ \answerBlank[5cm]
\answerExplain{$\frac{3}{8}, \; 0.39, \; 40\%$}{Convert: $40\% = 0.4$, $\frac{3}{8} = 0.375$. Order: $0.375 < 0.39 < 0.4$.}

\prob Compare $-\dfrac{2}{3}$ and $-65\%$. \answerBlank[3cm]
\answerExplain{$-\frac{2}{3} < -65\%$}{$-\frac{2}{3} \approx -66.7\%$ and $-65\% = -0.65$. Since $-0.667 < -0.65$, the fraction is less.}

\end{practiceBox}

% ── Practice D: True or False ────────────────────────────────────────────────
\begin{practiceBox}[funPurple]{True or False?}

\trueOrFalse{$-\frac{1}{2} > -\frac{1}{3}$}
\answerExplain{False}{$-\frac{1}{2} = -0.5$ and $-\frac{1}{3} \approx -0.333$. Since $-0.5 < -0.333$, the statement is false.}

\trueOrFalse{The absolute value of a number can be negative.}
\answerExplain{False}{Absolute value measures distance, which is always $\geq 0$.}

\trueOrFalse{$|{-5}| = |5|$}
\answerExplain{True}{Both are $5$ units from $0$: $|{-5}| = 5$ and $|5| = 5$.}

\trueOrFalse{$0.25 > \frac{1}{3}$}
\answerExplain{False}{$\frac{1}{3} \approx 0.333$, and $0.25 < 0.333$.}

\end{practiceBox}

% ── Word Problems ────────────────────────────────────────────────────────────
\begin{practiceBox}[funBlue]{\faGlobe~Word Problems}

\wordProblem{A submarine is at $-150$ feet. It rises $150$ feet. What is its new depth? Express as an integer.}{feet}
\answerExplain{$0$ feet (sea level)}{$-150 + 150 = 0$. The submarine is at sea level because $-150$ and $150$ are opposites.}

\wordProblem{Three cities recorded overnight temperatures: City~A at $-\frac{3}{2}°$C, City~B at $-1.8°$C, and City~C at $-\frac{7}{4}°$C. Rank them from coldest to warmest.}{city order}
\answerExplain{City~C, City~B, City~A}{Convert: A $= -1.5°$C, B $= -1.8°$C, C $= -1.75°$C. Order: $-1.8 < -1.75 < -1.5$, so C is coldest and A is warmest.}

\wordProblem{Mount Everest is $8{,}849$~m above sea level. The Mariana Trench is $10{,}994$~m below. Write each as an integer and compare their absolute values.}{meters}
\answerExplain{Everest: $+8{,}849$; Trench: $-10{,}994$; Trench has greater absolute value}{$|{-10{,}994}| = 10{,}994 > 8{,}849 = |8{,}849|$. The trench is farther from sea level.}

\end{practiceBox}

% ── Challenge ────────────────────────────────────────────────────────────────
\begin{challengeBox}
\prob If $|x| = 7$, what are the possible values of $x$? \answerBlank[3cm]
\answerExplain{$x = 7$ or $x = -7$}{Absolute value $= 7$ means the number is $7$ units from $0$: either $7$ or $-7$.}

\prob Place these numbers on a number line and identify the two that are closest to each other: $-\frac{5}{8}, \; -0.6, \; \frac{1}{4}, \; 0.2, \; -\frac{3}{10}$ \answerBlank[5cm]
\answerExplain{$-\frac{5}{8}$ and $-0.6$ are closest (distance $0.025$)}{Decimals: $-0.625, -0.6, 0.25, 0.2, -0.3$. Differences: $|{-0.625 - (-0.6)}| = 0.025$, $|0.25 - 0.2| = 0.05$. The pair $-\frac{5}{8}$ and $-0.6$ are only $0.025$ apart.}
\end{challengeBox}

\encouragement{You've mastered comparing and ordering rational numbers in every form --- fractions, decimals, and percents!}
